%%%%%%%%%%%%%%%%%%%%%%%%%%%%%%%%%%%%%%%%%%%%%%%%%%%%%%%%%%%%
% 4.4 GEOMETRIC TOPOLOGY
% The Composition of Advanced Mathematics
%%%%%%%%%%%%%%%%%%%%%%%%%%%%%%%%%%%%%%%%%%%%%%%%%%%%%%%%%%%%

\section{Geometric Topology}
\label{sec:geometric-topology}

\prerequisite{
Point-Set Topology, Algebraic Topology, Differential Topology,
introductory Group Theory, and basic Geometry.
}

\learningobjectives{
By the end of this section, the reader should be able to:
\begin{itemize}
    \item explain the goals of geometric topology;
    \item distinguish geometric topology from differential geometry;
    \item describe embeddings and ambient isotopy;
    \item explain why dimension plays a central role in topology;
    \item classify compact connected surfaces conceptually;
    \item work with orientable and nonorientable surfaces;
    \item compute Euler characteristic for standard surfaces;
    \item describe connected sums;
    \item explain triangulations and simplicial decompositions;
    \item describe handle decompositions;
    \item explain surgery as a topological operation;
    \item describe fundamental groups of surfaces;
    \item recognize lens spaces and other basic \(3\)-manifolds;
    \item explain Heegaard splittings;
    \item describe incompressible surfaces;
    \item explain prime decomposition of \(3\)-manifolds;
    \item state the Sphere Theorem conceptually;
    \item state the Loop Theorem conceptually;
    \item explain Dehn surgery;
    \item explain hyperbolic structures on \(3\)-manifolds;
    \item state the role of geometrization;
    \item explain the Poincar\'e Conjecture;
    \item describe high-dimensional surgery theory;
    \item distinguish tame and wild embeddings;
    \item connect geometric topology with knot theory,
    manifold theory, algebraic topology, and geometry.
\end{itemize}
}

%%%%%%%%%%%%%%%%%%%%%%%%%%%%%%%%%%%%%%%%%%%%%%%%%%%%%%%%%%%%
\subsection{What Is Geometric Topology?}
\label{subsec:what-is-geometric-topology}
%%%%%%%%%%%%%%%%%%%%%%%%%%%%%%%%%%%%%%%%%%%%%%%%%%%%%%%%%%%%

Geometric topology studies manifolds and embeddings using geometric,
combinatorial, and topological techniques.

A central goal is to understand the possible shapes of spaces while
retaining only those geometric features that are topologically meaningful.

Typical questions include:

\begin{itemize}
    \item Which manifolds are homeomorphic?
    \item How can manifolds be decomposed into simpler pieces?
    \item How can one manifold be constructed from another?
    \item How can curves and surfaces sit inside higher-dimensional spaces?
    \item Which geometric structures can a topological manifold support?
\end{itemize}

\begin{conceptbox}
Geometric topology studies topology through concrete geometric objects such
as manifolds, surfaces, knots, links, triangulations, handles, and
decompositions.
\end{conceptbox}

%%%%%%%%%%%%%%%%%%%%%%%%%%%%%%%%%%%%%%%%%%%%%%%%%%%%%%%%%%%%
\subsection{Topology Versus Geometry}
\label{subsec:topology-versus-geometry-geometric}
%%%%%%%%%%%%%%%%%%%%%%%%%%%%%%%%%%%%%%%%%%%%%%%%%%%%%%%%%%%%

A geometric object may carry many different kinds of information.

For example, a surface may have:

\[
\boxed{
\text{lengths},
\quad
\text{angles},
\quad
\text{curvature},
\quad
\text{topology}.
}
\]

Topological properties survive homeomorphisms.

Metric properties generally do not.

For example, stretching a sphere into an ellipsoid changes distances and
curvatures but not its topological type.

Thus

\[
\boxed{
S^2
\cong
\text{ellipsoid}
}
\]

as topological spaces.

%%%%%%%%%%%%%%%%%%%%%%%%%%%%%%%%%%%%%%%%%%%%%%%%%%%%%%%%%%%%
\subsection{Geometric Topology Versus Differential Geometry}
\label{subsec:geometric-topology-vs-differential-geometry}
%%%%%%%%%%%%%%%%%%%%%%%%%%%%%%%%%%%%%%%%%%%%%%%%%%%%%%%%%%%%

Differential geometry typically studies structures such as

\[
\boxed{
\text{metrics},
\quad
\text{connections},
\quad
\text{curvature},
\quad
\text{geodesics}.
}
\]

Geometric topology instead emphasizes

\[
\boxed{
\text{homeomorphism type},
\quad
\text{isotopy},
\quad
\text{decomposition},
\quad
\text{embedding},
\quad
\text{global structure}.
}
\]

The two subjects strongly interact.

For example, hyperbolic geometry became fundamental to the topology of
\(3\)-manifolds.

%%%%%%%%%%%%%%%%%%%%%%%%%%%%%%%%%%%%%%%%%%%%%%%%%%%%%%%%%%%%
\subsection{The Role of Dimension}
\label{subsec:dimension-geometric-topology}
%%%%%%%%%%%%%%%%%%%%%%%%%%%%%%%%%%%%%%%%%%%%%%%%%%%%%%%%%%%%

Topology changes dramatically with dimension.

One often organizes manifold topology by dimension:

\[
\boxed{
1,\quad2,\quad3,\quad4,\quad n\geq5.
}
\]

Different dimensions require different techniques.

Dimension \(2\) is governed by surface classification.

Dimension \(3\) is closely connected with knots and geometric structures.

Dimension \(4\) exhibits exceptional smooth and topological behavior.

Dimensions \(5\) and higher admit powerful surgery and cobordism methods.

%%%%%%%%%%%%%%%%%%%%%%%%%%%%%%%%%%%%%%%%%%%%%%%%%%%%%%%%%%%%
\subsection{One-Dimensional Manifolds}
\label{subsec:one-dimensional-manifolds}
%%%%%%%%%%%%%%%%%%%%%%%%%%%%%%%%%%%%%%%%%%%%%%%%%%%%%%%%%%%%

Connected \(1\)-manifolds are especially simple.

Under standard Hausdorff and second-countability assumptions, a connected
\(1\)-manifold without boundary is homeomorphic to either

\[
\boxed{
S^1
\qquad\text{or}\qquad
\R.
}
\]

Compact connected \(1\)-manifolds without boundary are therefore
homeomorphic to

\[
\boxed{
S^1.
}
\]

Compact connected \(1\)-manifolds with nonempty boundary are homeomorphic to

\[
\boxed{
[0,1].
}
\]

%%%%%%%%%%%%%%%%%%%%%%%%%%%%%%%%%%%%%%%%%%%%%%%%%%%%%%%%%%%%
\subsection{Surfaces}
\label{subsec:geometric-topology-surfaces}
%%%%%%%%%%%%%%%%%%%%%%%%%%%%%%%%%%%%%%%%%%%%%%%%%%%%%%%%%%%%

A \emph{surface} is a \(2\)-dimensional manifold.

Locally, every point of a surface resembles

\[
\R^2
\]

or, for a surface with boundary, a half-plane.

Globally, however, surfaces can have very different topology.

Examples include

\[
\boxed{
S^2,
\quad
T^2,
\quad
\mathbb{RP}^2,
\quad
\text{Klein bottle}.
}
\]

%%%%%%%%%%%%%%%%%%%%%%%%%%%%%%%%%%%%%%%%%%%%%%%%%%%%%%%%%%%%
\subsection{Orientable Surfaces}
\label{subsec:orientable-surfaces-geometric}
%%%%%%%%%%%%%%%%%%%%%%%%%%%%%%%%%%%%%%%%%%%%%%%%%%%%%%%%%%%%

A connected orientable closed surface is determined by its genus.

The genus

\[
\boxed{
g
}
\]

counts the number of handles.

For example,

\[
g=0
\]

gives the sphere,

\[
g=1
\]

gives the torus, and

\[
g=2
\]

gives the double torus.

We commonly denote the closed orientable surface of genus \(g\) by

\[
\boxed{
\Sigma_g.
}
\]

The uppercase Greek letter

\[
\Sigma
\]

is read ``Sigma.''

%%%%%%%%%%%%%%%%%%%%%%%%%%%%%%%%%%%%%%%%%%%%%%%%%%%%%%%%%%%%
\subsection{Euler Characteristic of Orientable Surfaces}
\label{subsec:euler-orientable-surfaces}
%%%%%%%%%%%%%%%%%%%%%%%%%%%%%%%%%%%%%%%%%%%%%%%%%%%%%%%%%%%%

For a closed connected orientable surface of genus \(g\),

\[
\boxed{
\chi(\Sigma_g)=2-2g.
}
\]

Thus

\[
\chi(S^2)=2,
\]

\[
\chi(T^2)=0,
\]

and

\[
\chi(\Sigma_2)=-2.
\]

%%%%%%%%%%%%%%%%%%%%%%%%%%%%%%%%%%%%%%%%%%%%%%%%%%%%%%%%%%%%
\subsection{Nonorientable Surfaces}
\label{subsec:nonorientable-surfaces-geometric}
%%%%%%%%%%%%%%%%%%%%%%%%%%%%%%%%%%%%%%%%%%%%%%%%%%%%%%%%%%%%

Closed connected nonorientable surfaces are constructed as connected sums
of copies of the real projective plane.

Write

\[
\boxed{
N_k
=
\underbrace{
\mathbb{RP}^2
\#\cdots\#
\mathbb{RP}^2
}_{k\text{ copies}}.
}
\]

The symbol

\[
\#
\]

is read ``connected sum.''

The integer \(k\) is sometimes called the nonorientable genus or
crosscap number.

%%%%%%%%%%%%%%%%%%%%%%%%%%%%%%%%%%%%%%%%%%%%%%%%%%%%%%%%%%%%
\subsection{Euler Characteristic of Nonorientable Surfaces}
\label{subsec:euler-nonorientable-surfaces}
%%%%%%%%%%%%%%%%%%%%%%%%%%%%%%%%%%%%%%%%%%%%%%%%%%%%%%%%%%%%

For

\[
N_k,
\]

the Euler characteristic is

\[
\boxed{
\chi(N_k)=2-k.
}
\]

For example,

\[
\chi(\mathbb{RP}^2)=1.
\]

The Klein bottle is homeomorphic to

\[
\mathbb{RP}^2\#\mathbb{RP}^2,
\]

so

\[
\boxed{
\chi(\text{Klein bottle})=0.
}
\]

%%%%%%%%%%%%%%%%%%%%%%%%%%%%%%%%%%%%%%%%%%%%%%%%%%%%%%%%%%%%
\subsection{Classification of Closed Surfaces}
\label{subsec:classification-closed-surfaces}
%%%%%%%%%%%%%%%%%%%%%%%%%%%%%%%%%%%%%%%%%%%%%%%%%%%%%%%%%%%%

One of the foundational classification theorems in topology completely
describes compact connected surfaces without boundary.

\begin{theorem}[Classification of Closed Surfaces]
Every compact connected surface without boundary is homeomorphic to exactly
one of the following types:

\[
\boxed{
S^2,
}
\]

\[
\boxed{
\Sigma_g
\qquad
(g\geq1),
}
\]

or

\[
\boxed{
N_k
\qquad
(k\geq1).
}
\]

Equivalently, every closed connected surface is either orientable of some
genus \(g\) or nonorientable of some crosscap number \(k\).
\end{theorem}

%%%%%%%%%%%%%%%%%%%%%%%%%%%%%%%%%%%%%%%%%%%%%%%%%%%%%%%%%%%%
\subsection{Why Surface Classification Is Remarkable}
\label{subsec:why-surface-classification}
%%%%%%%%%%%%%%%%%%%%%%%%%%%%%%%%%%%%%%%%%%%%%%%%%%%%%%%%%%%%

A surface may initially be presented in many ways:

\[
\boxed{
\text{polygon identifications},
\quad
\text{triangulations},
\quad
\text{equations},
\quad
\text{embedded surfaces}.
}
\]

The classification theorem says that all closed connected surfaces reduce
to a short canonical list.

Thus two such surfaces can be distinguished by information such as

\[
\boxed{
\text{orientability}
\quad\text{and}\quad
\chi.
}
\]

%%%%%%%%%%%%%%%%%%%%%%%%%%%%%%%%%%%%%%%%%%%%%%%%%%%%%%%%%%%%
\subsection{Surfaces with Boundary}
\label{subsec:surfaces-with-boundary-geometric}
%%%%%%%%%%%%%%%%%%%%%%%%%%%%%%%%%%%%%%%%%%%%%%%%%%%%%%%%%%%%

Compact surfaces may also contain boundary components.

An orientable compact connected surface can be described by

\[
\boxed{
(g,b),
}
\]

where

\[
g=\text{genus},
\qquad
b=\text{number of boundary components}.
\]

Its Euler characteristic is

\[
\boxed{
\chi=2-2g-b.
}
\]

For a nonorientable compact connected surface with crosscap number \(k\),

\[
\boxed{
\chi=2-k-b.
}
\]

%%%%%%%%%%%%%%%%%%%%%%%%%%%%%%%%%%%%%%%%%%%%%%%%%%%%%%%%%%%%
\subsection{Connected Sum}
\label{subsec:connected-sum-geometric}
%%%%%%%%%%%%%%%%%%%%%%%%%%%%%%%%%%%%%%%%%%%%%%%%%%%%%%%%%%%%

\begin{definitionbox}{Connected Sum}
Let \(M\) and \(N\) be connected \(n\)-manifolds.

Remove the interiors of small embedded \(n\)-balls from each manifold and
identify the resulting boundary spheres.

The resulting manifold is called the \emph{connected sum} and is denoted

\[
\boxed{
M\#N.
}
\]
\end{definitionbox}

For surfaces, connected sum provides a natural way to add handles or
crosscaps.

%%%%%%%%%%%%%%%%%%%%%%%%%%%%%%%%%%%%%%%%%%%%%%%%%%%%%%%%%%%%
\subsection{Connected Sum and Euler Characteristic}
\label{subsec:connected-sum-euler}
%%%%%%%%%%%%%%%%%%%%%%%%%%%%%%%%%%%%%%%%%%%%%%%%%%%%%%%%%%%%

For closed connected surfaces,

\[
\boxed{
\chi(M\#N)
=
\chi(M)+\chi(N)-2.
}
\]

For example,

\[
T^2\#T^2
\]

has Euler characteristic

\[
0+0-2=-2.
\]

Therefore it is the orientable genus-\(2\) surface.

%%%%%%%%%%%%%%%%%%%%%%%%%%%%%%%%%%%%%%%%%%%%%%%%%%%%%%%%%%%%
\subsection{Worked Example: Identifying an Orientable Surface}
\label{subsec:worked-identifying-surface}
%%%%%%%%%%%%%%%%%%%%%%%%%%%%%%%%%%%%%%%%%%%%%%%%%%%%%%%%%%%%

\begin{workedexamplebox}{Finding the Genus}

Suppose \(M\) is a closed connected orientable surface with

\[
\chi(M)=-6.
\]

For an orientable closed surface,

\[
\chi(M)=2-2g.
\]

Therefore

\[
-6=2-2g.
\]

Hence

\[
-8=-2g,
\]

so

\[
g=4.
\]

\begin{answerbox}
\[
\boxed{
M\cong\Sigma_4.
}
\]
\end{answerbox}
\end{workedexamplebox}

%%%%%%%%%%%%%%%%%%%%%%%%%%%%%%%%%%%%%%%%%%%%%%%%%%%%%%%%%%%%
\subsection{Polygonal Representations of Surfaces}
\label{subsec:polygonal-representations}
%%%%%%%%%%%%%%%%%%%%%%%%%%%%%%%%%%%%%%%%%%%%%%%%%%%%%%%%%%%%

Many surfaces can be represented by polygons whose edges are identified.

For example, the torus has the standard edge word

\[
\boxed{
aba^{-1}b^{-1}.
}
\]

The projective plane has edge word

\[
\boxed{
aa.
}
\]

The Klein bottle may be represented by

\[
\boxed{
aba^{-1}b.
}
\]

These representations encode global topology through boundary
identifications.

%%%%%%%%%%%%%%%%%%%%%%%%%%%%%%%%%%%%%%%%%%%%%%%%%%%%%%%%%%%%
\subsection{Fundamental Polygon of an Orientable Surface}
\label{subsec:fundamental-polygon-orientable}
%%%%%%%%%%%%%%%%%%%%%%%%%%%%%%%%%%%%%%%%%%%%%%%%%%%%%%%%%%%%

The closed orientable genus-\(g\) surface admits a polygonal presentation

\[
\boxed{
a_1b_1a_1^{-1}b_1^{-1}
\cdots
a_gb_ga_g^{-1}b_g^{-1}.
}
\]

Each pair

\[
a_i,b_i
\]

corresponds geometrically to one handle.

%%%%%%%%%%%%%%%%%%%%%%%%%%%%%%%%%%%%%%%%%%%%%%%%%%%%%%%%%%%%
\subsection{Fundamental Group of an Orientable Surface}
\label{subsec:fundamental-group-orientable-surface}
%%%%%%%%%%%%%%%%%%%%%%%%%%%%%%%%%%%%%%%%%%%%%%%%%%%%%%%%%%%%

The fundamental group of

\[
\Sigma_g
\]

has presentation

\[
\boxed{
\pi_1(\Sigma_g)
=
\left\langle
a_1,b_1,\ldots,a_g,b_g
\;\middle|\;
[a_1,b_1]\cdots[a_g,b_g]=1
\right\rangle.
}
\]

Here

\[
[a,b]
=
aba^{-1}b^{-1}
\]

is the group commutator.

For the torus,

\[
\boxed{
\pi_1(T^2)
\cong
\Z^2.
}
\]

%%%%%%%%%%%%%%%%%%%%%%%%%%%%%%%%%%%%%%%%%%%%%%%%%%%%%%%%%%%%
\subsection{Fundamental Group of a Nonorientable Surface}
\label{subsec:fundamental-group-nonorientable-surface}
%%%%%%%%%%%%%%%%%%%%%%%%%%%%%%%%%%%%%%%%%%%%%%%%%%%%%%%%%%%%

For the closed nonorientable surface \(N_k\),

\[
\boxed{
\pi_1(N_k)
=
\left\langle
a_1,\ldots,a_k
\;\middle|\;
a_1^2a_2^2\cdots a_k^2=1
\right\rangle.
}
\]

This illustrates how geometric decomposition produces algebraic invariants.

%%%%%%%%%%%%%%%%%%%%%%%%%%%%%%%%%%%%%%%%%%%%%%%%%%%%%%%%%%%%
\subsection{Triangulations}
\label{subsec:triangulations-geometric-topology}
%%%%%%%%%%%%%%%%%%%%%%%%%%%%%%%%%%%%%%%%%%%%%%%%%%%%%%%%%%%%

A triangulation decomposes a space into simplices.

For a surface, these simplices are primarily

\[
\boxed{
\text{vertices},
\quad
\text{edges},
\quad
\text{triangles}.
}
\]

If a finite triangulation has

\[
V
\]

vertices,

\[
E
\]

edges, and

\[
F
\]

faces, then

\[
\boxed{
\chi=V-E+F.
}
\]

%%%%%%%%%%%%%%%%%%%%%%%%%%%%%%%%%%%%%%%%%%%%%%%%%%%%%%%%%%%%
\subsection{Why Triangulations Matter}
\label{subsec:why-triangulations-matter}
%%%%%%%%%%%%%%%%%%%%%%%%%%%%%%%%%%%%%%%%%%%%%%%%%%%%%%%%%%%%

Triangulations convert continuous spaces into finite combinatorial data.

This makes it possible to use

\[
\boxed{
\text{combinatorics}
\longrightarrow
\text{algebra}
\longrightarrow
\text{topology}.
}
\]

Triangulations are especially important in

\[
\boxed{
\text{simplicial homology},
\quad
\text{PL topology},
\quad
\text{computational topology}.
}
\]

%%%%%%%%%%%%%%%%%%%%%%%%%%%%%%%%%%%%%%%%%%%%%%%%%%%%%%%%%%%%
\subsection{Piecewise-Linear Topology}
\label{subsec:pl-topology}
%%%%%%%%%%%%%%%%%%%%%%%%%%%%%%%%%%%%%%%%%%%%%%%%%%%%%%%%%%%%

Piecewise-linear topology, often abbreviated

\[
\boxed{
\mathrm{PL}
}
\]

topology, studies spaces and maps that are linear on the simplices of
suitable triangulations.

It provides an intermediate category between smooth and purely topological
manifolds.

One encounters the hierarchy

\[
\boxed{
\text{smooth}
\longrightarrow
\mathrm{PL}
\longrightarrow
\text{topological}.
}
\]

The relationships among these categories depend strongly on dimension.

%%%%%%%%%%%%%%%%%%%%%%%%%%%%%%%%%%%%%%%%%%%%%%%%%%%%%%%%%%%%
\subsection{Handle Decompositions}
\label{subsec:handle-decompositions-geometric}
%%%%%%%%%%%%%%%%%%%%%%%%%%%%%%%%%%%%%%%%%%%%%%%%%%%%%%%%%%%%

Differential topology introduced handles through Morse theory.

An \(n\)-dimensional \(k\)-handle is

\[
\boxed{
D^k\times D^{n-k}.
}
\]

A manifold can often be constructed by attaching handles successively.

Thus

\[
\boxed{
\text{manifold}
=
\text{collection of handles}
+
\text{attaching information}.
}
\]

%%%%%%%%%%%%%%%%%%%%%%%%%%%%%%%%%%%%%%%%%%%%%%%%%%%%%%%%%%%%
\subsection{Handles on Surfaces}
\label{subsec:handles-on-surfaces}
%%%%%%%%%%%%%%%%%%%%%%%%%%%%%%%%%%%%%%%%%%%%%%%%%%%%%%%%%%%%

For a \(2\)-manifold, the possible handle indices are

\[
\boxed{
0,\quad1,\quad2.
}
\]

A \(0\)-handle is a disk.

A \(1\)-handle is

\[
D^1\times D^1.
\]

A \(2\)-handle is another disk attached along its boundary.

For a torus, one may use

\[
\boxed{
1\text{ zero-handle},
\quad
2\text{ one-handles},
\quad
1\text{ two-handle}.
}
\]

This gives

\[
\chi(T^2)=1-2+1=0.
\]

%%%%%%%%%%%%%%%%%%%%%%%%%%%%%%%%%%%%%%%%%%%%%%%%%%%%%%%%%%%%
\subsection{Euler Characteristic from Handles}
\label{subsec:euler-from-handles}
%%%%%%%%%%%%%%%%%%%%%%%%%%%%%%%%%%%%%%%%%%%%%%%%%%%%%%%%%%%%

If a finite handle decomposition contains \(h_k\) handles of index \(k\),
then

\[
\boxed{
\chi(M)
=
\sum_k(-1)^kh_k.
}
\]

This parallels the corresponding formula for CW complexes and Morse
critical points.

%%%%%%%%%%%%%%%%%%%%%%%%%%%%%%%%%%%%%%%%%%%%%%%%%%%%%%%%%%%%
\subsection{Surgery}
\label{subsec:topological-surgery}
%%%%%%%%%%%%%%%%%%%%%%%%%%%%%%%%%%%%%%%%%%%%%%%%%%%%%%%%%%%%

Surgery modifies a manifold by removing one standard piece and replacing it
with another having the same boundary.

In an \(n\)-manifold, a \(k\)-surgery replaces

\[
\boxed{
S^k\times D^{n-k}
}
\]

with

\[
\boxed{
D^{k+1}\times S^{n-k-1}.
}
\]

Their boundaries agree:

\[
\boxed{
S^k\times S^{n-k-1}.
}
\]

%%%%%%%%%%%%%%%%%%%%%%%%%%%%%%%%%%%%%%%%%%%%%%%%%%%%%%%%%%%%
\subsection{Geometric Meaning of Surgery}
\label{subsec:geometric-meaning-surgery}
%%%%%%%%%%%%%%%%%%%%%%%%%%%%%%%%%%%%%%%%%%%%%%%%%%%%%%%%%%%%

Surgery can

\[
\boxed{
\text{create tunnels},
\quad
\text{remove tunnels},
\quad
\text{change connectivity},
\quad
\text{alter homotopy groups}.
}
\]

It is one of the principal construction and classification techniques in
high-dimensional topology.

%%%%%%%%%%%%%%%%%%%%%%%%%%%%%%%%%%%%%%%%%%%%%%%%%%%%%%%%%%%%
\subsection{Embeddings}
\label{subsec:geometric-topology-embeddings}
%%%%%%%%%%%%%%%%%%%%%%%%%%%%%%%%%%%%%%%%%%%%%%%%%%%%%%%%%%%%

An embedding

\[
f:M\hookrightarrow N
\]

places one topological space inside another without identifying distinct
points or changing its intrinsic topology.

Geometric topology asks not only whether an embedding exists, but how many
essentially different embeddings exist.

For example,

\[
\boxed{
S^1\hookrightarrow S^3
}
\]

leads directly to knot theory.

%%%%%%%%%%%%%%%%%%%%%%%%%%%%%%%%%%%%%%%%%%%%%%%%%%%%%%%%%%%%
\subsection{Ambient Isotopy}
\label{subsec:geometric-topology-ambient-isotopy}
%%%%%%%%%%%%%%%%%%%%%%%%%%%%%%%%%%%%%%%%%%%%%%%%%%%%%%%%%%%%

Two embedded objects are considered geometrically equivalent when one can
be continuously moved to the other through the surrounding space without
tearing or self-intersection.

For smooth manifolds, this is formalized by ambient isotopy.

Thus geometric topology often classifies

\[
\boxed{
\text{embeddings modulo ambient isotopy}.
}
\]

%%%%%%%%%%%%%%%%%%%%%%%%%%%%%%%%%%%%%%%%%%%%%%%%%%%%%%%%%%%%
\subsection{Tame and Wild Embeddings}
\label{subsec:tame-wild-embeddings}
%%%%%%%%%%%%%%%%%%%%%%%%%%%%%%%%%%%%%%%%%%%%%%%%%%%%%%%%%%%%

Not every topological embedding behaves like a smooth or polygonal
embedding.

A \emph{tame} embedding is, roughly, one equivalent to a well-behaved
standard embedding such as a PL or smooth one.

A \emph{wild} embedding has pathological local behavior.

%%%%%%%%%%%%%%%%%%%%%%%%%%%%%%%%%%%%%%%%%%%%%%%%%%%%%%%%%%%%
\subsection{Alexander Horned Sphere}
\label{subsec:alexander-horned-sphere}
%%%%%%%%%%%%%%%%%%%%%%%%%%%%%%%%%%%%%%%%%%%%%%%%%%%%%%%%%%%%

A famous example of a wild embedding is the Alexander horned sphere.

It is an embedding

\[
\boxed{
S^2\hookrightarrow\R^3
}
\]

whose image is topologically a sphere but is embedded in an extremely
complicated way.

Its complement does not behave like the complement of the standard round
sphere.

This demonstrates that

\[
\boxed{
\text{intrinsic topology}
\neq
\text{embedding topology}.
}
\]

%%%%%%%%%%%%%%%%%%%%%%%%%%%%%%%%%%%%%%%%%%%%%%%%%%%%%%%%%%%%
\subsection{Three-Manifolds}
\label{subsec:three-manifolds-geometric}
%%%%%%%%%%%%%%%%%%%%%%%%%%%%%%%%%%%%%%%%%%%%%%%%%%%%%%%%%%%%

A \(3\)-manifold is a space locally modeled on

\[
\R^3.
\]

Examples include

\[
\boxed{
S^3,
\quad
T^3,
\quad
S^2\times S^1,
\quad
\mathbb{RP}^3.
}
\]

Three-dimensional topology is exceptionally rich because surfaces can sit
inside \(3\)-manifolds and divide them into meaningful pieces.

%%%%%%%%%%%%%%%%%%%%%%%%%%%%%%%%%%%%%%%%%%%%%%%%%%%%%%%%%%%%
\subsection{The Three-Sphere}
\label{subsec:three-sphere-geometric}
%%%%%%%%%%%%%%%%%%%%%%%%%%%%%%%%%%%%%%%%%%%%%%%%%%%%%%%%%%%%

The \(3\)-sphere is

\[
\boxed{
S^3
=
\{
(x_1,x_2,x_3,x_4)\in\R^4:
x_1^2+x_2^2+x_3^2+x_4^2=1
\}.
}
\]

Although \(S^3\) cannot be embedded in \(\R^3\), it is the natural ambient
space for classical knot theory.

Adding one point at infinity to \(\R^3\) gives

\[
\boxed{
\R^3\cup\{\infty\}
\cong
S^3.
}
\]

%%%%%%%%%%%%%%%%%%%%%%%%%%%%%%%%%%%%%%%%%%%%%%%%%%%%%%%%%%%%
\subsection{Lens Spaces}
\label{subsec:lens-spaces}
%%%%%%%%%%%%%%%%%%%%%%%%%%%%%%%%%%%%%%%%%%%%%%%%%%%%%%%%%%%%

Lens spaces provide an important family of closed \(3\)-manifolds.

They can be constructed by gluing two solid tori along their boundaries
using an appropriate homeomorphism.

They are traditionally denoted

\[
\boxed{
L(p,q).
}
\]

For suitable coprime integers \(p\) and \(q\), their fundamental groups
satisfy

\[
\boxed{
\pi_1(L(p,q))
\cong
\Z/p\Z.
}
\]

Different lens spaces provide early examples showing that manifold
classification requires subtle invariants.

%%%%%%%%%%%%%%%%%%%%%%%%%%%%%%%%%%%%%%%%%%%%%%%%%%%%%%%%%%%%
\subsection{Solid Tori}
\label{subsec:solid-tori}
%%%%%%%%%%%%%%%%%%%%%%%%%%%%%%%%%%%%%%%%%%%%%%%%%%%%%%%%%%%%

A solid torus is

\[
\boxed{
S^1\times D^2.
}
\]

Its boundary is

\[
\boxed{
\partial(S^1\times D^2)
=
S^1\times S^1
=
T^2.
}
\]

Solid tori appear repeatedly in

\[
\boxed{
\text{knot complements},
\quad
\text{Heegaard splittings},
\quad
\text{Dehn surgery}.
}
\]

%%%%%%%%%%%%%%%%%%%%%%%%%%%%%%%%%%%%%%%%%%%%%%%%%%%%%%%%%%%%
\subsection{Heegaard Splittings}
\label{subsec:heegaard-splittings}
%%%%%%%%%%%%%%%%%%%%%%%%%%%%%%%%%%%%%%%%%%%%%%%%%%%%%%%%%%%%

A closed connected orientable \(3\)-manifold can be decomposed into two
handlebodies glued along their common boundary.

Schematically,

\[
\boxed{
M
=
H_g\cup_{\varphi}H_g.
}
\]

Here

\[
H_g
\]

is a genus-\(g\) handlebody and

\[
\varphi:\partial H_g\to\partial H_g
\]

is a gluing homeomorphism.

This is called a \emph{Heegaard splitting}.

%%%%%%%%%%%%%%%%%%%%%%%%%%%%%%%%%%%%%%%%%%%%%%%%%%%%%%%%%%%%
\subsection{Heegaard Genus}
\label{subsec:heegaard-genus}
%%%%%%%%%%%%%%%%%%%%%%%%%%%%%%%%%%%%%%%%%%%%%%%%%%%%%%%%%%%%

The smallest genus among all Heegaard splittings of \(M\) is called the
\emph{Heegaard genus}.

It provides an important measure of the complexity of a \(3\)-manifold.

For example,

\[
S^3
\]

admits a genus-\(0\) splitting into two \(3\)-balls.

Lens spaces admit genus-\(1\) splittings into two solid tori.

%%%%%%%%%%%%%%%%%%%%%%%%%%%%%%%%%%%%%%%%%%%%%%%%%%%%%%%%%%%%
\subsection{Compressing Disks}
\label{subsec:compressing-disks}
%%%%%%%%%%%%%%%%%%%%%%%%%%%%%%%%%%%%%%%%%%%%%%%%%%%%%%%%%%%%

Let

\[
S\subset M
\]

be an embedded surface.

A compressing disk is an embedded disk

\[
D\subset M
\]

whose boundary lies in \(S\) and represents an essential simple closed curve
in \(S\).

Compressing along such a disk can simplify the topology of the surface.

%%%%%%%%%%%%%%%%%%%%%%%%%%%%%%%%%%%%%%%%%%%%%%%%%%%%%%%%%%%%
\subsection{Incompressible Surfaces}
\label{subsec:incompressible-surfaces}
%%%%%%%%%%%%%%%%%%%%%%%%%%%%%%%%%%%%%%%%%%%%%%%%%%%%%%%%%%%%

An embedded surface is called \emph{incompressible} when it has no
appropriate compressing disk, subject to standard exceptions and boundary
conditions.

Algebraically, incompressibility is closely related to injectivity of the
induced map

\[
\boxed{
\pi_1(S)\to\pi_1(M).
}
\]

Incompressible surfaces are central tools for decomposing \(3\)-manifolds.

%%%%%%%%%%%%%%%%%%%%%%%%%%%%%%%%%%%%%%%%%%%%%%%%%%%%%%%%%%%%
\subsection{Irreducible Three-Manifolds}
\label{subsec:irreducible-three-manifolds}
%%%%%%%%%%%%%%%%%%%%%%%%%%%%%%%%%%%%%%%%%%%%%%%%%%%%%%%%%%%%

\begin{definitionbox}{Irreducible \(3\)-Manifold}
A connected \(3\)-manifold is \emph{irreducible} if every embedded
\(2\)-sphere bounds a \(3\)-ball.
\end{definitionbox}

The condition says that embedded spheres do not reveal nontrivial connected
sum decompositions.

%%%%%%%%%%%%%%%%%%%%%%%%%%%%%%%%%%%%%%%%%%%%%%%%%%%%%%%%%%%%
\subsection{Prime Three-Manifolds}
\label{subsec:prime-three-manifolds}
%%%%%%%%%%%%%%%%%%%%%%%%%%%%%%%%%%%%%%%%%%%%%%%%%%%%%%%%%%%%

\begin{definitionbox}{Prime \(3\)-Manifold}
A connected \(3\)-manifold \(M\) is \emph{prime} if

\[
M\cong M_1\#M_2
\]

implies that one of the summands is homeomorphic to \(S^3\).
\end{definitionbox}

Prime manifolds play a role analogous to prime numbers.

They are the basic building blocks for connected sums.

%%%%%%%%%%%%%%%%%%%%%%%%%%%%%%%%%%%%%%%%%%%%%%%%%%%%%%%%%%%%
\subsection{Prime Decomposition Theorem}
\label{subsec:prime-decomposition-three-manifolds}
%%%%%%%%%%%%%%%%%%%%%%%%%%%%%%%%%%%%%%%%%%%%%%%%%%%%%%%%%%%%

A foundational theorem states that a compact connected orientable
\(3\)-manifold can be decomposed as a connected sum of prime
\(3\)-manifolds.

Under the standard hypotheses, this decomposition is unique up to order and
homeomorphism of the prime factors.

Schematically,

\[
\boxed{
M
\cong
M_1\#M_2\#\cdots\#M_r.
}
\]

This is the \(3\)-manifold analogue of integer prime factorization.

%%%%%%%%%%%%%%%%%%%%%%%%%%%%%%%%%%%%%%%%%%%%%%%%%%%%%%%%%%%%
\subsection{Sphere Theorem}
\label{subsec:sphere-theorem}
%%%%%%%%%%%%%%%%%%%%%%%%%%%%%%%%%%%%%%%%%%%%%%%%%%%%%%%%%%%%

The Sphere Theorem connects homotopy information with embedded spheres.

Conceptually, it states that under appropriate hypotheses, nontrivial
second homotopy

\[
\boxed{
\pi_2(M)\neq0
}
\]

can be represented by a suitably nontrivial embedded \(2\)-sphere.

This converts an abstract homotopy-theoretic obstruction into a geometric
surface.

%%%%%%%%%%%%%%%%%%%%%%%%%%%%%%%%%%%%%%%%%%%%%%%%%%%%%%%%%%%%
\subsection{Loop Theorem}
\label{subsec:loop-theorem}
%%%%%%%%%%%%%%%%%%%%%%%%%%%%%%%%%%%%%%%%%%%%%%%%%%%%%%%%%%%%

The Loop Theorem similarly converts algebraic information about fundamental
groups into embedded disks.

Roughly, when a boundary loop becomes trivial in the ambient manifold in a
nontrivial way, one can often find an embedded disk realizing that
compression.

The theorem is fundamental to the study of incompressible surfaces.

%%%%%%%%%%%%%%%%%%%%%%%%%%%%%%%%%%%%%%%%%%%%%%%%%%%%%%%%%%%%
\subsection{Dehn Filling}
\label{subsec:dehn-filling}
%%%%%%%%%%%%%%%%%%%%%%%%%%%%%%%%%%%%%%%%%%%%%%%%%%%%%%%%%%%%

Suppose a \(3\)-manifold \(M\) has torus boundary

\[
\partial M\cong T^2.
\]

A \emph{Dehn filling} attaches a solid torus

\[
S^1\times D^2
\]

to this boundary.

The attaching map is determined by specifying which slope on the boundary
torus becomes the meridian of the new solid torus.

%%%%%%%%%%%%%%%%%%%%%%%%%%%%%%%%%%%%%%%%%%%%%%%%%%%%%%%%%%%%
\subsection{Slopes on a Torus}
\label{subsec:slopes-torus}
%%%%%%%%%%%%%%%%%%%%%%%%%%%%%%%%%%%%%%%%%%%%%%%%%%%%%%%%%%%%

Choose a meridian-longitude basis

\[
\boxed{
(\mu,\lambda)
}
\]

for

\[
H_1(T^2;\Z).
\]

The lowercase Greek letters

\[
\mu
\qquad\text{and}\qquad
\lambda
\]

are read ``mu'' and ``lambda.''

A primitive homology class

\[
\boxed{
p\mu+q\lambda
}
\]

determines a slope, where

\[
\gcd(p,q)=1.
\]

This slope determines a Dehn filling.

%%%%%%%%%%%%%%%%%%%%%%%%%%%%%%%%%%%%%%%%%%%%%%%%%%%%%%%%%%%%
\subsection{Dehn Surgery on Knots}
\label{subsec:dehn-surgery-knots}
%%%%%%%%%%%%%%%%%%%%%%%%%%%%%%%%%%%%%%%%%%%%%%%%%%%%%%%%%%%%

Let

\[
K\subset S^3
\]

be a knot.

Remove a tubular neighborhood

\[
N(K).
\]

The complement

\[
\boxed{
S^3\setminus\operatorname{int}N(K)
}
\]

has torus boundary.

Attaching a solid torus back using a different boundary identification
produces a new \(3\)-manifold.

This process is called \emph{Dehn surgery} on the knot.

%%%%%%%%%%%%%%%%%%%%%%%%%%%%%%%%%%%%%%%%%%%%%%%%%%%%%%%%%%%%
\subsection{Surgery Description of Three-Manifolds}
\label{subsec:surgery-description-three-manifolds}
%%%%%%%%%%%%%%%%%%%%%%%%%%%%%%%%%%%%%%%%%%%%%%%%%%%%%%%%%%%%

A remarkable result states that every closed connected orientable
\(3\)-manifold can be obtained by surgery on a link in \(S^3\).

Thus the classification of \(3\)-manifolds is deeply connected with the
study of knots and links.

Schematically,

\[
\boxed{
\text{link in }S^3
\overset{\text{surgery}}{\longrightarrow}
\text{closed orientable }3\text{-manifold}.
}
\]

%%%%%%%%%%%%%%%%%%%%%%%%%%%%%%%%%%%%%%%%%%%%%%%%%%%%%%%%%%%%
\subsection{Geometric Structures on Three-Manifolds}
\label{subsec:geometric-structures-three-manifolds}
%%%%%%%%%%%%%%%%%%%%%%%%%%%%%%%%%%%%%%%%%%%%%%%%%%%%%%%%%%%%

A major theme in \(3\)-manifold topology is the interaction between topology
and geometry.

One asks whether a manifold can be equipped locally with a standard
homogeneous geometry.

The most familiar possibilities include

\[
\boxed{
\text{spherical},
\quad
\text{Euclidean},
\quad
\text{hyperbolic}.
}
\]

Three-dimensional topology requires additional model geometries as well.

%%%%%%%%%%%%%%%%%%%%%%%%%%%%%%%%%%%%%%%%%%%%%%%%%%%%%%%%%%%%
\subsection{Thurston's Eight Geometries}
\label{subsec:thurston-eight-geometries}
%%%%%%%%%%%%%%%%%%%%%%%%%%%%%%%%%%%%%%%%%%%%%%%%%%%%%%%%%%%%

The eight model geometries appearing in Thurston's geometrization program
are

\[
\boxed{
S^3,
\quad
\R^3,
\quad
\mathbb{H}^3,
}
\]

\[
\boxed{
S^2\times\R,
\quad
\mathbb{H}^2\times\R,
}
\]

together with the geometries commonly denoted

\[
\boxed{
\widetilde{\mathrm{SL}_2(\R)},
\quad
\mathrm{Nil},
\quad
\mathrm{Sol}.
}
\]

These geometries provide canonical local models for pieces of
\(3\)-manifolds.

%%%%%%%%%%%%%%%%%%%%%%%%%%%%%%%%%%%%%%%%%%%%%%%%%%%%%%%%%%%%
\subsection{Hyperbolic Three-Manifolds}
\label{subsec:hyperbolic-three-manifolds}
%%%%%%%%%%%%%%%%%%%%%%%%%%%%%%%%%%%%%%%%%%%%%%%%%%%%%%%%%%%%

A hyperbolic \(3\)-manifold is locally modeled on hyperbolic space

\[
\boxed{
\mathbb{H}^3.
}
\]

Its metric has constant sectional curvature

\[
\boxed{
K=-1.
}
\]

Many knot and link complements admit complete finite-volume hyperbolic
metrics.

This creates a powerful bridge between

\[
\boxed{
\text{knot topology}
\quad\text{and}\quad
\text{hyperbolic geometry}.
}
\]

%%%%%%%%%%%%%%%%%%%%%%%%%%%%%%%%%%%%%%%%%%%%%%%%%%%%%%%%%%%%
\subsection{Mostow Rigidity}
\label{subsec:mostow-rigidity}
%%%%%%%%%%%%%%%%%%%%%%%%%%%%%%%%%%%%%%%%%%%%%%%%%%%%%%%%%%%%

In dimensions at least \(3\), finite-volume hyperbolic geometry is remarkably
rigid.

For closed hyperbolic manifolds, and more generally in the appropriate
finite-volume setting, the fundamental group determines the hyperbolic
structure up to isometry.

This is a consequence of Mostow--Prasad rigidity.

Thus geometric quantities such as hyperbolic volume become topological
invariants.

%%%%%%%%%%%%%%%%%%%%%%%%%%%%%%%%%%%%%%%%%%%%%%%%%%%%%%%%%%%%
\subsection{Hyperbolic Volume}
\label{subsec:hyperbolic-volume-topology}
%%%%%%%%%%%%%%%%%%%%%%%%%%%%%%%%%%%%%%%%%%%%%%%%%%%%%%%%%%%%

For a finite-volume hyperbolic \(3\)-manifold \(M\), its hyperbolic volume

\[
\boxed{
\operatorname{Vol}(M)
}
\]

is determined by its topology.

This is unusual because volume is normally a geometric quantity rather than
a topological one.

Rigidity turns geometry into topology.

%%%%%%%%%%%%%%%%%%%%%%%%%%%%%%%%%%%%%%%%%%%%%%%%%%%%%%%%%%%%
\subsection{JSJ Decomposition}
\label{subsec:jsj-decomposition}
%%%%%%%%%%%%%%%%%%%%%%%%%%%%%%%%%%%%%%%%%%%%%%%%%%%%%%%%%%%%

After prime decomposition, many irreducible \(3\)-manifolds can be further
cut along incompressible tori.

The resulting pieces are organized by the
Jaco--Shalen--Johannson decomposition, usually abbreviated

\[
\boxed{
\mathrm{JSJ}.
}
\]

This decomposition separates the manifold into pieces that are, broadly,
Seifert-fibered or hyperbolic in the geometrization picture.

%%%%%%%%%%%%%%%%%%%%%%%%%%%%%%%%%%%%%%%%%%%%%%%%%%%%%%%%%%%%
\subsection{Geometrization}
\label{subsec:geometrization}
%%%%%%%%%%%%%%%%%%%%%%%%%%%%%%%%%%%%%%%%%%%%%%%%%%%%%%%%%%%%

Thurston's Geometrization Conjecture proposed that compact
\(3\)-manifolds can be decomposed into canonical pieces, each admitting one
of the eight model geometries.

The conjecture was proved through the Ricci-flow program developed by
Hamilton and completed by Perelman.

Its proof transformed the classification theory of \(3\)-manifolds.

%%%%%%%%%%%%%%%%%%%%%%%%%%%%%%%%%%%%%%%%%%%%%%%%%%%%%%%%%%%%
\subsection{Poincar\'e Conjecture}
\label{subsec:poincare-conjecture}
%%%%%%%%%%%%%%%%%%%%%%%%%%%%%%%%%%%%%%%%%%%%%%%%%%%%%%%%%%%%

The Poincar\'e Conjecture is a special case of geometrization.

\begin{theorem}[Poincar\'e Conjecture]
Every closed simply connected topological \(3\)-manifold is homeomorphic to

\[
\boxed{
S^3.
}
\]
\end{theorem}

Thus, in dimension three,

\[
\boxed{
\pi_1(M)=0
}
\]

for a closed \(3\)-manifold implies

\[
\boxed{
M\cong S^3.
}
\]

%%%%%%%%%%%%%%%%%%%%%%%%%%%%%%%%%%%%%%%%%%%%%%%%%%%%%%%%%%%%
\subsection{Why the Poincar\'e Conjecture Matters}
\label{subsec:why-poincare-matters}
%%%%%%%%%%%%%%%%%%%%%%%%%%%%%%%%%%%%%%%%%%%%%%%%%%%%%%%%%%%%

For surfaces, simply connected closed manifolds are easy to classify.

The corresponding question in dimension three resisted proof for roughly a
century.

Its resolution illustrates a central feature of geometric topology:

\[
\boxed{
\text{topological classification}
+
\text{geometric structures}
+
\text{analysis}
}
\]

can combine to solve problems inaccessible to any one viewpoint alone.

%%%%%%%%%%%%%%%%%%%%%%%%%%%%%%%%%%%%%%%%%%%%%%%%%%%%%%%%%%%%
\subsection{Ricci Flow and Topology}
\label{subsec:ricci-flow-topology}
%%%%%%%%%%%%%%%%%%%%%%%%%%%%%%%%%%%%%%%%%%%%%%%%%%%%%%%%%%%%

Ricci flow evolves a Riemannian metric according to

\[
\boxed{
\frac{\partial g}{\partial t}
=
-2\operatorname{Ric}(g).
}
\]

Here

\[
\operatorname{Ric}
\]

denotes Ricci curvature.

The flow attempts to regularize the geometry of a manifold.

Singularities may form, requiring careful analysis and topological
modification.

The full theory belongs to differential geometry and geometric analysis,
but its consequences are fundamental to geometric topology.

%%%%%%%%%%%%%%%%%%%%%%%%%%%%%%%%%%%%%%%%%%%%%%%%%%%%%%%%%%%%
\subsection{Four-Manifolds}
\label{subsec:four-manifolds-geometric-topology}
%%%%%%%%%%%%%%%%%%%%%%%%%%%%%%%%%%%%%%%%%%%%%%%%%%%%%%%%%%%%

Dimension four behaves differently from both low-dimensional and
high-dimensional topology.

Topological \(4\)-manifolds can exhibit phenomena such as

\[
\boxed{
\text{exotic smooth structures}
}
\]

and especially subtle relationships between topology and differentiability.

As seen in Differential Topology,

\[
\R^4
\]

admits exotic smooth structures.

No other Euclidean space

\[
\R^n
\]

has this particular phenomenon.

%%%%%%%%%%%%%%%%%%%%%%%%%%%%%%%%%%%%%%%%%%%%%%%%%%%%%%%%%%%%
\subsection{Intersection Forms of Four-Manifolds}
\label{subsec:intersection-forms-four-manifolds}
%%%%%%%%%%%%%%%%%%%%%%%%%%%%%%%%%%%%%%%%%%%%%%%%%%%%%%%%%%%%

For a closed oriented \(4\)-manifold, intersections of \(2\)-dimensional
homology classes define a bilinear form

\[
\boxed{
Q_M:
H_2(M;\Z)/\text{torsion}
\times
H_2(M;\Z)/\text{torsion}
\to
\Z.
}
\]

This is called the \emph{intersection form}.

It is a central invariant in \(4\)-manifold topology.

The special role of dimension four occurs partly because

\[
\boxed{
2+2=4,
}
\]

so surfaces generically intersect in isolated points.

%%%%%%%%%%%%%%%%%%%%%%%%%%%%%%%%%%%%%%%%%%%%%%%%%%%%%%%%%%%%
\subsection{High-Dimensional Manifolds}
\label{subsec:high-dimensional-manifolds-geometric}
%%%%%%%%%%%%%%%%%%%%%%%%%%%%%%%%%%%%%%%%%%%%%%%%%%%%%%%%%%%%

In dimensions at least five, techniques such as

\[
\boxed{
\text{handle theory},
\quad
\text{surgery theory},
\quad
\text{cobordism},
\quad
\text{homotopy theory}
}
\]

become especially powerful.

Many classification questions can be translated into algebraic conditions.

This does not make high-dimensional topology easy, but it gives the subject
a systematic framework unavailable in some lower dimensions.

%%%%%%%%%%%%%%%%%%%%%%%%%%%%%%%%%%%%%%%%%%%%%%%%%%%%%%%%%%%%
\subsection{The \texorpdfstring{\(h\)}{h}-Cobordism Theorem}
\label{subsec:h-cobordism-theorem}
%%%%%%%%%%%%%%%%%%%%%%%%%%%%%%%%%%%%%%%%%%%%%%%%%%%%%%%%%%%%

An \(h\)-cobordism is a cobordism

\[
W
\]

between manifolds \(M_0\) and \(M_1\) such that the inclusions

\[
M_0\hookrightarrow W
\qquad\text{and}\qquad
M_1\hookrightarrow W
\]

are homotopy equivalences.

A foundational high-dimensional theorem states that, under suitable
dimension and simple-connectivity assumptions, an \(h\)-cobordism is
diffeomorphic to a product.

Conceptually,

\[
\boxed{
W\cong M\times[0,1].
}
\]

%%%%%%%%%%%%%%%%%%%%%%%%%%%%%%%%%%%%%%%%%%%%%%%%%%%%%%%%%%%%
\subsection{The \texorpdfstring{\(s\)}{s}-Cobordism Theorem}
\label{subsec:s-cobordism}
%%%%%%%%%%%%%%%%%%%%%%%%%%%%%%%%%%%%%%%%%%%%%%%%%%%%%%%%%%%%

When the fundamental group is nontrivial, an additional invariant called
Whitehead torsion appears.

The \(s\)-Cobordism Theorem states, under the appropriate
high-dimensional hypotheses, that an \(h\)-cobordism is a product exactly
when its Whitehead torsion vanishes.

Thus

\[
\boxed{
\text{geometric product structure}
\longleftrightarrow
\text{algebraic torsion obstruction}.
}
\]

%%%%%%%%%%%%%%%%%%%%%%%%%%%%%%%%%%%%%%%%%%%%%%%%%%%%%%%%%%%%
\subsection{Surgery Theory in High Dimensions}
\label{subsec:surgery-theory-high-dimensions}
%%%%%%%%%%%%%%%%%%%%%%%%%%%%%%%%%%%%%%%%%%%%%%%%%%%%%%%%%%%%

High-dimensional surgery theory attempts to modify maps and manifolds until
a homotopy equivalence can be improved to a homeomorphism or
diffeomorphism.

The general strategy is

\[
\boxed{
\text{homotopy data}
\longrightarrow
\text{surgery operations}
\longrightarrow
\text{manifold classification}.
}
\]

Obstructions to completing the process are encoded algebraically in
surgery obstruction groups.

%%%%%%%%%%%%%%%%%%%%%%%%%%%%%%%%%%%%%%%%%%%%%%%%%%%%%%%%%%%%
\subsection{Local Versus Global Topology}
\label{subsec:local-global-geometric-topology}
%%%%%%%%%%%%%%%%%%%%%%%%%%%%%%%%%%%%%%%%%%%%%%%%%%%%%%%%%%%%

Every \(n\)-manifold locally resembles

\[
\R^n.
\]

Therefore manifolds are locally indistinguishable at the purely topological
level.

Their differences are global.

For example,

\[
S^2
\qquad\text{and}\qquad
T^2
\]

are both locally modeled on

\[
\R^2,
\]

but globally

\[
\boxed{
\pi_1(S^2)=0
}
\]

while

\[
\boxed{
\pi_1(T^2)\cong\Z^2.
}
\]

%%%%%%%%%%%%%%%%%%%%%%%%%%%%%%%%%%%%%%%%%%%%%%%%%%%%%%%%%%%%
\subsection{Geometric Decomposition as a General Strategy}
\label{subsec:geometric-decomposition-strategy}
%%%%%%%%%%%%%%%%%%%%%%%%%%%%%%%%%%%%%%%%%%%%%%%%%%%%%%%%%%%%

A recurring strategy throughout geometric topology is:

\[
\boxed{
\text{cut a complicated space into simpler pieces}.
}
\]

Examples include

\[
\boxed{
\text{surface decomposition},
}
\]

\[
\boxed{
\text{prime decomposition},
}
\]

\[
\boxed{
\text{JSJ decomposition},
}
\]

\[
\boxed{
\text{handle decomposition}.
}
\]

After understanding the pieces, one studies how they are glued together.

%%%%%%%%%%%%%%%%%%%%%%%%%%%%%%%%%%%%%%%%%%%%%%%%%%%%%%%%%%%%
\subsection{Algebraic Invariants in Geometric Topology}
\label{subsec:algebraic-invariants-geometric-topology}
%%%%%%%%%%%%%%%%%%%%%%%%%%%%%%%%%%%%%%%%%%%%%%%%%%%%%%%%%%%%

Geometric topology repeatedly uses algebraic invariants to distinguish
spaces.

Important examples include

\[
\boxed{
\pi_1(M),
}
\]

\[
\boxed{
H_k(M),
}
\]

\[
\boxed{
H^k(M),
}
\]

\[
\boxed{
\chi(M).
}
\]

In higher-dimensional topology one also encounters

\[
\boxed{
\text{intersection forms},
\quad
\text{characteristic classes},
\quad
\text{torsion invariants}.
}
\]

%%%%%%%%%%%%%%%%%%%%%%%%%%%%%%%%%%%%%%%%%%%%%%%%%%%%%%%%%%%%
\subsection{Worked Example: Euler Characteristic of a Surface}
\label{subsec:worked-euler-surface-geometric}
%%%%%%%%%%%%%%%%%%%%%%%%%%%%%%%%%%%%%%%%%%%%%%%%%%%%%%%%%%%%

\begin{workedexamplebox}{Classifying a Nonorientable Surface}

Suppose \(M\) is a closed connected nonorientable surface with

\[
\chi(M)=-3.
\]

For a nonorientable surface,

\[
\chi(N_k)=2-k.
\]

Therefore

\[
-3=2-k.
\]

Hence

\[
k=5.
\]

\begin{answerbox}
\[
\boxed{
M\cong N_5
=
\#^5\mathbb{RP}^2.
}
\]
\end{answerbox}
\end{workedexamplebox}

%%%%%%%%%%%%%%%%%%%%%%%%%%%%%%%%%%%%%%%%%%%%%%%%%%%%%%%%%%%%
\subsection{Worked Example: Dimension of an Intersection}
\label{subsec:worked-intersection-geometric}
%%%%%%%%%%%%%%%%%%%%%%%%%%%%%%%%%%%%%%%%%%%%%%%%%%%%%%%%%%%%

\begin{workedexamplebox}{Surfaces in a Three-Manifold}

Suppose two embedded surfaces \(A\) and \(B\) in a \(3\)-manifold \(M\)
intersect transversely.

Then

\[
\dim A=2,
\qquad
\dim B=2,
\qquad
\dim M=3.
\]

Therefore

\[
\dim(A\cap B)
=
2+2-3
=
1.
\]

\begin{answerbox}
\[
\boxed{
A\cap B
\text{ is locally a }1\text{-dimensional submanifold}.
}
\]
\end{answerbox}

Thus generic transverse surfaces in a \(3\)-manifold intersect along curves.
\end{workedexamplebox}

%%%%%%%%%%%%%%%%%%%%%%%%%%%%%%%%%%%%%%%%%%%%%%%%%%%%%%%%%%%%
\subsection{Worked Example: Connected Sum}
\label{subsec:worked-connected-sum}
%%%%%%%%%%%%%%%%%%%%%%%%%%%%%%%%%%%%%%%%%%%%%%%%%%%%%%%%%%%%

\begin{workedexamplebox}{Euler Characteristic of a Triple Torus}

Consider

\[
M
=
T^2\#T^2\#T^2.
\]

Since each torus has

\[
\chi(T^2)=0,
\]

we calculate

\[
\chi(T^2\#T^2)
=
0+0-2
=
-2.
\]

Then

\[
\chi(M)
=
-2+0-2
=
-4.
\]

Alternatively, \(M\) has genus \(3\), so

\[
\chi(M)
=
2-2(3)
=
-4.
\]

\begin{answerbox}
\[
\boxed{
\chi(T^2\#T^2\#T^2)=-4.
}
\]
\end{answerbox}
\end{workedexamplebox}

%%%%%%%%%%%%%%%%%%%%%%%%%%%%%%%%%%%%%%%%%%%%%%%%%%%%%%%%%%%%
\subsection{Common Errors}
\label{subsec:geometric-topology-common-errors}
%%%%%%%%%%%%%%%%%%%%%%%%%%%%%%%%%%%%%%%%%%%%%%%%%%%%%%%%%%%%

\subsubsection{Confusing Geometric Topology with Differential Geometry}

Geometric topology does not primarily study metric quantities such as
curvature.

It studies topological structures using geometric methods.

%%%%%%%%%%%%%%%%%%%%%%%%%%%%%%%%%%%%%%%%%%%%%%%%%%%%%%%%%%%%
\subsubsection{Assuming Locally Similar Manifolds Are Globally Equivalent}

Every \(n\)-manifold locally resembles \(\R^n\).

This does not imply that all \(n\)-manifolds are homeomorphic.

%%%%%%%%%%%%%%%%%%%%%%%%%%%%%%%%%%%%%%%%%%%%%%%%%%%%%%%%%%%%
\subsubsection{Confusing Genus with Euler Characteristic}

For a closed orientable surface,

\[
\chi=2-2g.
\]

The quantities \(g\) and \(\chi\) are related but are not identical.

%%%%%%%%%%%%%%%%%%%%%%%%%%%%%%%%%%%%%%%%%%%%%%%%%%%%%%%%%%%%
\subsubsection{Using the Orientable Formula for Nonorientable Surfaces}

For orientable closed surfaces,

\[
\chi=2-2g.
\]

For nonorientable closed surfaces,

\[
\chi=2-k.
\]

%%%%%%%%%%%%%%%%%%%%%%%%%%%%%%%%%%%%%%%%%%%%%%%%%%%%%%%%%%%%
\subsubsection{Assuming the Torus Is Simply Connected}

The torus has

\[
\boxed{
\pi_1(T^2)\cong\Z^2.
}
\]

%%%%%%%%%%%%%%%%%%%%%%%%%%%%%%%%%%%%%%%%%%%%%%%%%%%%%%%%%%%%
\subsubsection{Confusing Connected Sum with Disjoint Union}

Connected sum joins manifolds after deleting balls.

Disjoint union keeps the spaces separate.

Thus

\[
M\#N
\]

and

\[
M\sqcup N
\]

are entirely different constructions.

%%%%%%%%%%%%%%%%%%%%%%%%%%%%%%%%%%%%%%%%%%%%%%%%%%%%%%%%%%%%
\subsubsection{Assuming Every Embedding Is Equivalent}

Two embeddings of the same abstract space can be topologically different.

Knots provide the most familiar examples.

%%%%%%%%%%%%%%%%%%%%%%%%%%%%%%%%%%%%%%%%%%%%%%%%%%%%%%%%%%%%
\subsubsection{Confusing Intrinsic and Embedding Topology}

The image of an embedded sphere may be intrinsically homeomorphic to
\(S^2\) while its complement has complicated topology.

The Alexander horned sphere illustrates this distinction.

%%%%%%%%%%%%%%%%%%%%%%%%%%%%%%%%%%%%%%%%%%%%%%%%%%%%%%%%%%%%
\subsubsection{Assuming Prime Means Irreducible}

For orientable \(3\)-manifolds, prime and irreducible are closely related
but are not identical notions.

In particular,

\[
S^2\times S^1
\]

is prime but not irreducible.

%%%%%%%%%%%%%%%%%%%%%%%%%%%%%%%%%%%%%%%%%%%%%%%%%%%%%%%%%%%%
\subsubsection{Assuming Dehn Surgery and General Surgery Are Identical}

Dehn surgery is a specialized \(3\)-dimensional construction, commonly
performed along knots or links.

General surgery theory applies in arbitrary dimensions.

%%%%%%%%%%%%%%%%%%%%%%%%%%%%%%%%%%%%%%%%%%%%%%%%%%%%%%%%%%%%
\subsubsection{Assuming All Three-Manifolds Are Hyperbolic}

Many \(3\)-manifolds are hyperbolic, but not all.

Geometrization requires several possible model geometries.

%%%%%%%%%%%%%%%%%%%%%%%%%%%%%%%%%%%%%%%%%%%%%%%%%%%%%%%%%%%%
\subsubsection{Assuming the Poincar\'e Conjecture Applies in Every Dimension}

The classical Poincar\'e Conjecture concerns closed simply connected
\(3\)-manifolds.

Analogous generalized Poincar\'e problems occur in other dimensions and
have different histories and methods.

%%%%%%%%%%%%%%%%%%%%%%%%%%%%%%%%%%%%%%%%%%%%%%%%%%%%%%%%%%%%
\subsubsection{Ignoring Dimension in Topological Arguments}

Methods that work in dimension \(2\), \(3\), or at least \(5\) may fail or
change substantially in dimension \(4\).

Dimension is structural, not merely a numerical parameter.

%%%%%%%%%%%%%%%%%%%%%%%%%%%%%%%%%%%%%%%%%%%%%%%%%%%%%%%%%%%%
\subsection{Applications and Connections}
\label{subsec:geometric-topology-applications}
%%%%%%%%%%%%%%%%%%%%%%%%%%%%%%%%%%%%%%%%%%%%%%%%%%%%%%%%%%%%

\begin{applicationbox}
\textbf{Algebraic Topology.}
Fundamental groups, homology, cohomology, and Euler characteristic provide
algebraic tools for distinguishing geometric spaces.

\medskip

\textbf{Differential Topology.}
Embeddings, isotopies, handles, Morse theory, transversality, and cobordism
supply smooth techniques for geometric-topological problems.

\medskip

\textbf{Differential Geometry.}
Hyperbolic geometry, curvature, and Ricci flow can reveal the topology of
manifolds.

\medskip

\textbf{Knot Theory.}
Knots and links are embeddings inside \(3\)-manifolds, and their complements
are fundamental objects of geometric topology.

\medskip

\textbf{Group Theory.}
Fundamental groups of manifolds encode geometric information. Geometric
group theory studies groups through spaces on which they act.

\medskip

\textbf{Complex Geometry.}
Many complex manifolds have underlying smooth and topological structures
whose classification requires geometric-topological techniques.

\medskip

\textbf{Physics.}
Manifold topology appears in relativity, gauge theories, field theories,
topological phases, and models of spacetime.

\medskip

\textbf{Computation.}
Triangulated manifolds allow algorithms to analyze topology using finite
combinatorial structures.
\end{applicationbox}

%%%%%%%%%%%%%%%%%%%%%%%%%%%%%%%%%%%%%%%%%%%%%%%%%%%%%%%%%%%%
\subsection{Exercises}
\label{subsec:geometric-topology-exercises}
%%%%%%%%%%%%%%%%%%%%%%%%%%%%%%%%%%%%%%%%%%%%%%%%%%%%%%%%%%%%

\begin{exercise}[1]
What is geometric topology?
\end{exercise}

\begin{exercise}[1]
Explain the difference between geometric topology and differential geometry.
\end{exercise}

\begin{exercise}[1]
What are the compact connected \(1\)-manifolds without boundary?
\end{exercise}

\begin{exercise}[1]
Define the genus of an orientable surface.
\end{exercise}

\begin{exercise}[1]
State the Euler characteristic formula for \(\Sigma_g\).
\end{exercise}

\begin{exercise}[1]
State the Euler characteristic formula for \(N_k\).
\end{exercise}

\begin{exercise}[1]
Define connected sum.
\end{exercise}

\begin{exercise}[1]
State the classification theorem for closed connected surfaces.
\end{exercise}

\begin{exercise}[1]
What is a triangulation?
\end{exercise}

\begin{exercise}[1]
What does \(\mathrm{PL}\) stand for?
\end{exercise}

\begin{exercise}[1]
What is a handle?
\end{exercise}

\begin{exercise}[1]
What is topological surgery?
\end{exercise}

\begin{exercise}[1]
Define an irreducible \(3\)-manifold.
\end{exercise}

\begin{exercise}[1]
Define a prime \(3\)-manifold.
\end{exercise}

\begin{exercise}[1]
What is a Heegaard splitting?
\end{exercise}

\begin{exercise}[1]
What is Dehn filling?
\end{exercise}

\begin{exercise}[1]
State the Poincar\'e Conjecture.
\end{exercise}

\begin{exercise}[2]
Calculate the Euler characteristic of the orientable surface of genus \(5\).
\end{exercise}

\begin{exercise}[2]
A closed orientable surface has Euler characteristic \(-8\).
Determine its genus.
\end{exercise}

\begin{exercise}[2]
Calculate the Euler characteristic of \(N_7\).
\end{exercise}

\begin{exercise}[2]
A closed nonorientable surface has Euler characteristic \(-4\).
Determine its crosscap number.
\end{exercise}

\begin{exercise}[2]
Calculate

\[
\chi(T^2\#T^2).
\]
\end{exercise}

\begin{exercise}[2]
Calculate

\[
\chi(\Sigma_3\#\Sigma_2).
\]
\end{exercise}

\begin{exercise}[2]
Determine the genus of

\[
\Sigma_3\#\Sigma_2.
\]
\end{exercise}

\begin{exercise}[2]
Compute the Euler characteristic of an orientable genus-\(3\) surface with
four boundary components.
\end{exercise}

\begin{exercise}[2]
Compute the Euler characteristic of a nonorientable surface with five
crosscaps and two boundary components.
\end{exercise}

\begin{exercise}[2]
Write the standard fundamental-group presentation for \(\Sigma_2\).
\end{exercise}

\begin{exercise}[2]
Explain why

\[
\pi_1(T^2)\cong\Z^2.
\]
\end{exercise}

\begin{exercise}[2]
If a triangulated surface has

\[
V=12,\qquad E=30,\qquad F=20,
\]

calculate its Euler characteristic.
\end{exercise}

\begin{exercise}[2]
If the surface in the previous exercise is closed and orientable, determine
its genus.
\end{exercise}

\begin{exercise}[2]
Explain why

\[
S^1\times D^2
\]

has torus boundary.
\end{exercise}

\begin{exercise}[2]
Describe the genus-\(0\) Heegaard splitting of \(S^3\).
\end{exercise}

\begin{exercise}[2]
Explain why lens spaces naturally have genus-\(1\) Heegaard splittings.
\end{exercise}

\begin{exercise}[2]
Explain geometrically why two transverse surfaces in a \(3\)-manifold
generically intersect along curves.
\end{exercise}

\begin{exercise}[3]
Derive

\[
\chi(M\#N)=\chi(M)+\chi(N)-2
\]

for closed surfaces.
\end{exercise}

\begin{exercise}[3]
Use connected sums to construct an orientable surface of genus \(4\).
\end{exercise}

\begin{exercise}[3]
Use connected sums to construct a nonorientable surface with crosscap number
\(6\).
\end{exercise}

\begin{exercise}[3]
Show that

\[
\mathbb{RP}^2\#\mathbb{RP}^2
\]

has Euler characteristic \(0\).
\end{exercise}

\begin{exercise}[3]
Explain how polygonal edge identifications produce the torus.
\end{exercise}

\begin{exercise}[3]
Derive the standard presentation

\[
\pi_1(\Sigma_g)
=
\left\langle
a_1,b_1,\ldots,a_g,b_g
\mid
[a_1,b_1]\cdots[a_g,b_g]=1
\right\rangle
\]

from its fundamental polygon.
\end{exercise}

\begin{exercise}[3]
Explain how a handle decomposition produces an Euler characteristic formula.
\end{exercise}

\begin{exercise}[3]
Describe geometrically what a \(1\)-surgery does to a suitable
\(3\)-manifold.
\end{exercise}

\begin{exercise}[3]
Explain why an incompressible surface should carry nontrivial information
about the fundamental group of the ambient manifold.
\end{exercise}

\begin{exercise}[3]
Explain why an irreducible \(3\)-manifold cannot contain an essential
embedded \(2\)-sphere.
\end{exercise}

\begin{exercise}[3]
Explain why

\[
S^2\times S^1
\]

is not irreducible.
\end{exercise}

\begin{exercise}[3]
Research why

\[
S^2\times S^1
\]

is nevertheless prime.
\end{exercise}

\begin{exercise}[3]
Describe the complement of the unknot in \(S^3\).
\end{exercise}

\begin{exercise}[3]
Explain how Dehn filling changes a manifold with torus boundary into a new
\(3\)-manifold.
\end{exercise}

\begin{exercise}[3]
Explain why a knot complement naturally has torus boundary.
\end{exercise}

\begin{exercise}[3]
Describe the relationship between Dehn surgery and lens spaces.
\end{exercise}

\begin{exercise}[3]
Explain why Mostow rigidity allows hyperbolic volume to become a topological
invariant.
\end{exercise}

\begin{exercise}[3]
Explain the role of incompressible tori in the JSJ decomposition.
\end{exercise}

\begin{exercise}[3]
Describe the relationship between prime decomposition, JSJ decomposition,
and geometrization.
\end{exercise}

\begin{exercise}[4]
Prove the classification of compact connected \(1\)-manifolds under the
standard manifold hypotheses.
\end{exercise}

\begin{exercise}[4]
Study a proof of the Classification Theorem for Compact Surfaces.
\end{exercise}

\begin{exercise}[4]
Research the relationship between triangulations and PL structures on
manifolds.
\end{exercise}

\begin{exercise}[4]
Study the Alexander horned sphere and explain why it is a wild embedding.
\end{exercise}

\begin{exercise}[4]
Research the Sphere Theorem and explain its significance in
\(3\)-manifold topology.
\end{exercise}

\begin{exercise}[4]
Research the Loop Theorem and explain its relationship with compressing
disks.
\end{exercise}

\begin{exercise}[4]
Study the proof of prime decomposition for compact orientable
\(3\)-manifolds.
\end{exercise}

\begin{exercise}[4]
Construct a Heegaard diagram for a lens space.
\end{exercise}

\begin{exercise}[4]
Research the Lickorish--Wallace theorem on surgery descriptions of closed
orientable \(3\)-manifolds.
\end{exercise}

\begin{exercise}[4]
Study Thurston's eight model geometries and give one example associated with
each geometry.
\end{exercise}

\begin{exercise}[4]
Research a finite-volume hyperbolic knot complement.
\end{exercise}

\begin{exercise}[4]
Explain how Mostow--Prasad rigidity connects fundamental groups with
hyperbolic geometry.
\end{exercise}

\begin{exercise}[4]
Research the JSJ decomposition and describe its role in geometrization.
\end{exercise}

\begin{exercise}[4]
Study the relationship between Ricci flow and the proof of geometrization.
\end{exercise}

\begin{exercise}[4]
Explain why the Poincar\'e Conjecture follows from geometrization.
\end{exercise}

\begin{exercise}[4]
Research intersection forms of closed oriented \(4\)-manifolds.
\end{exercise}

\begin{exercise}[4]
Research the \(h\)-Cobordism Theorem and identify its dimension hypotheses.
\end{exercise}

\begin{exercise}[4]
Research Whitehead torsion and the \(s\)-Cobordism Theorem.
\end{exercise}

\begin{exercise}[4]
Investigate the basic strategy of high-dimensional surgery theory.
\end{exercise}

%%%%%%%%%%%%%%%%%%%%%%%%%%%%%%%%%%%%%%%%%%%%%%%%%%%%%%%%%%%%
\subsection{Conceptual Summary}
\label{subsec:geometric-topology-summary}
%%%%%%%%%%%%%%%%%%%%%%%%%%%%%%%%%%%%%%%%%%%%%%%%%%%%%%%%%%%%

Geometric topology studies manifolds and embeddings through concrete
geometric and topological constructions.

Dimension strongly affects the behavior of manifolds:

\[
\boxed{
1,\quad2,\quad3,\quad4,\quad n\geq5
}
\]

each exhibit distinct phenomena.

Closed connected surfaces admit a complete classification.

Orientable surfaces satisfy

\[
\boxed{
\chi(\Sigma_g)=2-2g,
}
\]

while nonorientable surfaces satisfy

\[
\boxed{
\chi(N_k)=2-k.
}
\]

Connected sum provides a fundamental construction:

\[
\boxed{
M\#N.
}
\]

Triangulations replace continuous spaces with combinatorial models.

Handle decompositions build manifolds from pieces

\[
\boxed{
D^k\times D^{n-k}.
}
\]

Surgery modifies manifolds by replacing

\[
S^k\times D^{n-k}
\]

with

\[
D^{k+1}\times S^{n-k-1}.
\]

Three-dimensional topology introduces powerful decomposition techniques.

Prime decomposition gives

\[
\boxed{
M
\cong
M_1\#\cdots\#M_r.
}
\]

JSJ decomposition cuts appropriate irreducible manifolds along essential
tori.

Geometrization then relates the resulting pieces to standard geometric
models.

The Poincar\'e Conjecture identifies the simply connected closed
\(3\)-manifold:

\[
\boxed{
\pi_1(M)=0
\Longrightarrow
M\cong S^3.
}
\]

Hyperbolic topology produces the striking phenomenon that geometric
quantities such as volume can become topological invariants.

Dimension four contains exceptional phenomena involving smooth structures
and intersection forms.

In sufficiently high dimensions, cobordism and surgery provide systematic
classification methods.

Throughout geometric topology, one repeatedly follows the strategy

\[
\boxed{
\text{decompose}
\longrightarrow
\text{understand the pieces}
\longrightarrow
\text{understand the gluing}.
}
\]

%%%%%%%%%%%%%%%%%%%%%%%%%%%%%%%%%%%%%%%%%%%%%%%%%%%%%%%%%%%%
\subsection{Section Connections}
\label{subsec:geometric-topology-connections}
%%%%%%%%%%%%%%%%%%%%%%%%%%%%%%%%%%%%%%%%%%%%%%%%%%%%%%%%%%%%

\chapterconnection{
Geometric Topology combines ideas from Point-Set Topology, Algebraic
Topology, Differential Topology, Geometry, and Group Theory. Surface
classification provides the first major example of a complete manifold
classification. Fundamental groups and homology convert geometric
decompositions into algebraic information. Differential Topology supplies
embeddings, isotopies, handles, Morse theory, and cobordisms. In dimension
three, prime decomposition, incompressible surfaces, Dehn surgery, JSJ
decomposition, and hyperbolic geometry lead to Thurston's geometrization
picture and the Poincare Conjecture. In dimension four, intersection forms
and exotic smooth structures reveal exceptional behavior. In higher
dimensions, surgery and cobordism become major classification tools. The
study of embeddings \(S^1\hookrightarrow S^3\) leads directly into Knot
Theory, while the systematic study of manifolds continues in the next
section.
}

%%%%%%%%%%%%%%%%%%%%%%%%%%%%%%%%%%%%%%%%%%%%%%%%%%%%%%%%%%%%
% SECTION SOURCE TRACKING
%%%%%%%%%%%%%%%%%%%%%%%%%%%%%%%%%%%%%%%%%%%%%%%%%%%%%%%%%%%%

%------------------------------------------------------------
% 4.4 Geometric Topology
%------------------------------------------------------------
%
% Source basis:
%   - Standard geometric topology, surface topology,
%     manifold topology, and introductory 3-manifold theory.
%
% Used for:
%   - geometric topology motivation;
%   - role of dimension;
%   - 1-manifolds;
%   - surfaces;
%   - orientable and nonorientable surfaces;
%   - genus and crosscap number;
%   - Euler characteristic;
%   - classification of compact surfaces;
%   - surfaces with boundary;
%   - connected sums;
%   - polygonal representations;
%   - surface fundamental groups;
%   - triangulations;
%   - PL topology;
%   - handle decompositions;
%   - surgery;
%   - embeddings;
%   - ambient isotopy;
%   - tame and wild embeddings;
%   - Alexander horned sphere;
%   - 3-manifolds;
%   - S^3;
%   - lens spaces;
%   - solid tori;
%   - Heegaard splittings;
%   - compressing disks;
%   - incompressible surfaces;
%   - irreducible and prime 3-manifolds;
%   - prime decomposition;
%   - Sphere and Loop Theorems;
%   - Dehn filling and surgery;
%   - geometric structures on 3-manifolds;
%   - Thurston's eight geometries;
%   - hyperbolic 3-manifolds;
%   - Mostow--Prasad rigidity;
%   - hyperbolic volume;
%   - JSJ decomposition;
%   - geometrization;
%   - Poincare Conjecture;
%   - Ricci flow connection;
%   - 4-manifolds;
%   - intersection forms;
%   - high-dimensional manifolds;
%   - h-cobordism;
%   - s-cobordism;
%   - surgery theory.
%
% Notes:
%   - Fundamental groups, homology, cohomology, and Euler
%     characteristic are developed canonically in
%     Algebraic Topology.
%   - Smooth embeddings, transversality, Morse theory,
%     and cobordism are developed in Differential Topology.
%   - Hyperbolic geometry is introduced in Non-Euclidean
%     Geometry.
%   - Ricci flow belongs canonically to Differential
%     Geometry and Geometric Analysis.
%   - Knot complements, Reidemeister moves, knot invariants,
%     and link theory are developed in Knot Theory.
%   - General manifold theory is developed in the dedicated
%     Manifolds section.
%   - Advanced 3-manifold and 4-manifold classification is
%     introduced structurally rather than proved here.
%
%%%%%%%%%%%%%%%%%%%%%%%%%%%%%%%%%%%%%%%%%%%%%%%%%%%%%%%%%%%%